\selectlanguage{english}
\chapter*{Abstract}
\addcontentsline{toc}{chapter}{Abstract}

\noindent\textbf{Title:} The influence of confidence function on the quality of game recommendations in the implicit ALS model

\vspace{0.8em}

This thesis investigates the impact of confidence function design on the quality of recommendations generated by the implicit ALS model. In the method proposed by Hu, Koren, and Volinsky, observations enter the model with weights determined as a function of the intensity of user--item interactions, yet this choice is often treated as an implicit assumption. The thesis examines how transformations of interaction intensity affect recommendation quality and compares their impact with that of the model class.

The study was conducted on real-world data from the Steam platform, comprising users' game libraries together with the recorded cumulative playtime for each title. Five confidence function variants were compared: linear, logarithmic, logarithmic with saturation, logarithmic with offset, and a constant function that ignores playtime entirely. These were evaluated alongside the Bayesian Personalized Ranking model and a popularity-based recommender, with hyperparameters tuned independently using Bayesian optimization. Recommendation quality was assessed using ranking accuracy metrics as well as beyond-accuracy measures: catalog coverage, intra-list diversity, and popularity bias. The study also covered result stability, statistical significance testing with effect size estimation, and sensitivity to the confidence function parameters.

The results indicate that, for the dataset studied, monotonic weighting of confidence according to playtime does not improve recommendation quality. The variant with the constant confidence function achieved the highest accuracy, outperforming the logarithmic variant by 6.2\% and the linear variant by 8.1\% in terms of \ndcg{}. At the same time, the gap separating the implicit ALS models from the BPR model proved considerably larger than the entire range of differences observed within the implicit ALS family, while the BPR model achieved more than twice the catalog coverage. The choice of the model class therefore determines primarily the trade-off between the different dimensions of recommendation quality.

The results suggest that, in the analyzed dataset, the mere fact of owning a game was a sufficiently strong preference signal to match the performance of transformations that incorporated playtime. This indicates that, for the transformation variants considered, playtime information could not be leveraged in a way that improved recommendation quality. However, this does not imply that playtime itself lacks informative value. These conclusions apply specifically to the definition of the signal used in this study, namely the cumulative playtime accumulated over the entire period of game ownership.

\vspace{1em}
\noindent\textbf{Keywords:} recommender systems, implicit feedback, matrix factorization, implicit ALS, confidence function, Steam platform
\selectlanguage{polish}